\begin{abstract}
\rev{배틀로얄 게임에서는 여러 팀이 싸우는 동안 경기 구역이 단계적으로 줄어든다.
본 연구는 다음 단계 탈락 위험, 곧 새 구역이 공개될 때 살아 있는 팀이 그 구역으로의 축소가 끝나기 전에 탈락할 확률을 추정한다.
기존 설계는 다른 값을 예측한다. 하나는 개인의 생존을 분류하고, 다른 하나는 상대위험으로 팀의 순위를 매기는데, 이 순위로는 단계마다 위험이 얼마나 높은지 알 수 없다. 그래서 어떤 설계 선택이 예측을 좌우하는지 분명하지 않다.
본 연구는 줄어드는 구역을 생태학의 서식지 손실로 보고, 이 위험을 단계 단위로(이산시간 hazard) 모델링하며 각 팀을 네 가지 서식지 요인, 곧 구역에 대한 위치, 구역에 들어가기 위해 부족한 이동량, 주변 경쟁 팀 수, 남은 팀원 수로 기술한다.
PUBG 13,948경기의 기록 약 3억 2천만 건에서 예측 시점 이전의 기록만 쓴 팀--단계 표본 817,766개를 만들고, 두 기존 설계를 같은 표본에 맞추어 적용한 뒤, 맵, 1인칭·3인칭 시점, 팀 규모별 12개 코호트에서 설계를 한 단계씩 바꾸었다.
주변 경쟁 팀이 많고 남은 팀원이 적을수록 모든 코호트와 강건성 점검에서 위험이 높았으나, 이동 부족과의 연관은 주로 게임이 일반 경기에 채워 넣는 AI 팀 때문이었다.
본 연구의 설계는 모든 코호트에서 두 기존 설계보다 AUC가 높았고(0.778 대 0.686--0.763), 단계가 갈수록 위험이 높아지는 양상을 재현하였다. 순위 모델은 이를 거의 평평하게 예측하였다.
순위 모델의 신경망을 그대로 두고 학습 목적만 다음 단계 탈락에 대한 교차엔트로피로 바꾸면, 같은 단계의 팀끼리 비교한 AUC가 0.022, 같은 경기의 팀끼리 비교한 AUC가 0.009 올랐다. 이 이득은 단계마다 평균 위험을 조정하는 것으로는 생기지 않는다.
신경망 구조를 바꾸어도 같은 경기 안의 판별은 의미 있게 나아지지 않았고, 생존 분류기와의 차이는 주로 서식지 요인에서 왔다.
각 학습 목적은 자기 과제에서 가장 좋았고, 예측할 값에 맞는 학습 목적을 고르는 일이 신경망 구조보다 중요하였다.}
\end{abstract}

% 주제어는 EasyChair 제출 양식에 입력(본문에서는 분량 때문에 생략):
% 시공간 텔레메트리, 이산시간 hazard, 데이터 누수, 생존분석, e스포츠 분석
